\subsection{Removing unramified order from the remaining rank-one problem}

The following reduction does not assert the complete-local Betti bound
(F13.5). It identifies a bounded-degree geometric problem sufficient for
that bound on the metabelian systems actually used in F13. The
complete reduction passed an independent root proof audit.

\begin{lemma}[F13c: bounded ramified order after an unramified twist]
Retain $X_0,D_0,U_0$ and $t_i$ from F13, and choose $e\geq0$
such that $p^e>\max_i t_i$. Let $f:X_1\to X_0$ be any connected
finite \'etale cover, with $U_1=f^{-1}(U_0)$. Every character
\[
 \chi\in H^1(U_1,\mathbf Q_p/\mathbf Z_p)
\]
can be written $\chi=\alpha+\psi|_{U_1}$, where
$p^e\alpha=0$ and $\psi\in H^1(X_1,\mathbf Q_p/\mathbf Z_p)$.
The same $e$ works for every such cover. The Swan conductor of
$\alpha$ at every exceptional component equals that of $\chi$.
\end{lemma}
\begin{proof}
First, $H^1(X_1,\mathbf Q_p/\mathbf Z_p)$ is divisible. Indeed,
proper henselian specialization identifies the finite \'etale cover
categories of $X_1$ and its reduced proper nodal closed fibre $D_1$
\cite[Lemma 58.9.1, tag 0A48]{StacksProperEtaleSpecialization}.
Thus their groups of finite continuous characters agree. To verify
divisibility without any assertion about the nonabelian fundamental
group, use the Artin--Schreier sequence on $D_1$. The group
$H^1(D_1,\mathcal O_{D_1})$ is a finite-dimensional $k$-vector space,
and $H^2(D_1,\mathcal O_{D_1})=0$. On this vector space, Frobenius
is additive and $p$-semilinear. The map $F-1$ is therefore a
homomorphism of vector algebraic groups with differential $-1$.
It is \'etale and its image is an open subgroup of a connected
vector group, hence is the whole group. Since $k$ is algebraically
closed, it is surjective on $k$-points. Artin--Schreier now gives
$H^2(D_1,\mathbf F_p)=0$.

The exact sequence
\[
 0\longrightarrow\mathbf F_p\longrightarrow
 \mathbf Q_p/\mathbf Z_p\xrightarrow{p}
 \mathbf Q_p/\mathbf Z_p\longrightarrow0
\]
shows that multiplication by $p$ is surjective on
$H^1(D_1,\mathbf Q_p/\mathbf Z_p)$. All of its elements are finite
characters, and the same is true on $X_1$. This proves the required
divisibility.

The F13b conductor bound applies to every rank-one character on
$U_1$, not only to the systems in the statement of F13.
By that bound and the Witt-restriction calculation (F10.6), the character
$p^e\chi$ has conductor zero along every exceptional component.
It extends to $X_1$ by purity. Choose an unramified character
$\psi$ whose $p^e$ multiple is this extension, using divisibility.
Then $\alpha=\chi-\psi|_{U_1}$ is killed by $p^e$. An unramified
twist does not change a positive logarithmic break and preserves
conductor zero; hence the Swan conductors are unchanged.
\end{proof}

\begin{proposition}[F13d: a bounded-degree sufficient criterion]
Keep the notation and $e$ of F13c. Suppose that there is a constant
$C$ with the following property. For every connected finite \'etale
cover $X_1\to X_0$ of degree $d$ and every connected cyclic cover
$V\to U_1$ of degree dividing $p^e$ corresponding to a character
whose Swan conductors satisfy the F13b bounds, the normalization
of the complete local surface underlying $X_1$ in $V$ admits an
SNC resolution $Y$ dominating $X_1$ such that
\[
 d\bigl(\pi_1^t(V;Y,D_Y)\bigr)\leq Cd.
 \tag{F13.6}
\]
Then there is a uniform rank-linear $H^1(U_0,-)$ bound for every
coefficient-$\ell\ne p$ system occurring in F13. Consequently
(F13.6), together with the already proved F12, would settle the
uniform finite-generation assertion for quotients of derived
length at most three.
\end{proposition}
\begin{proof}
First consider a rank-one character $\theta$ on $U_1$ with finite
values in a field of characteristic $\ell\ne p$. Split its finite
image into its $p$-primary and prime-to-$p$ parts. By F13c, after
a finite extension of the coefficient field if necessary, write
\[
                         \theta=\alpha\otimes\beta,
\]
where $\alpha$ has order dividing $p^e$ and $\beta$ is tame
along $D_1$. Here the $p$-primary factor removed from $\alpha$
extends over $X_1$, while every prime-to-$p$ rank-one character
is tame. Let $V\to U_1$ be the connected cyclic cover associated
to $\alpha$. Its degree is a power of $p$, hence invertible in
the coefficient field. Inflation--restriction gives an injection
\[
                H^1(U_1,\theta)\longrightarrow H^1(V,\beta|_V).
 \tag{F13.7}
\]

The pullback $\beta|_V$ is tame along $D_Y$. For its unramified
$p$-primary factor this follows by pulling back the cover of
$X_1$ along the proper morphism $Y\to X_1$; for its prime-to-$p$
factor it follows from its prime-to-$p$ image at each DVR.
Every degree-one cocycle for this rank-one system also factors
through the tame quotient. Indeed, the corresponding affine
representation has image in the semidirect product of the additive
coefficient group with the image of $\beta$. On wild inertia
$\beta$ is trivial, so the cocycle restricts to a continuous
homomorphism from a pro-$p$ group to an additive $\ell$-group,
which is zero. The affine cover is therefore tame. Thus (F13.6)
and values of a cocycle on generators give
\[
          \dim H^1(V,\beta|_V)\leq Cd.
 \tag{F13.8}
\]

Now let $M$ be a system occurring in F13. Extend scalars to a
finite field which splits its finite monodromy group. Each simple
constituent $W$ is, by F13a, induced from a rank-one character
$\theta$ on a finite \'etale resolution cover $X_1\to X_0$,
of degree $d=\operatorname{rank}W$. Shapiro's lemma and
(F13.7)--(F13.8) give
\[
 \dim H^1(U_0,W)=\dim H^1(U_1,\theta)
                   \leq C\operatorname{rank}W.
\]
The long exact sequence for a composition series adds these
bounds, proving $\dim H^1(U_0,M)\leq C\operatorname{rank}M$.
Finite extension of the coefficient field preserves these
dimensions; for a fixed finite local system this follows by
tensoring the continuous cochain complexes, or by their filtered
finite-quotient description. Finally apply the fixed-cover
inflation--restriction and complement count already explained
after (F13.5).
\end{proof}

\paragraph{Scope of the reduction.}
The exponent $e$ in F13c is proved uniform. The generator bound
(F13.6) is a separate, unresolved assertion. Bounded degree of a
ramified cover alone has not been shown here to bound the genera
and incidence graph of a resolution linearly in $d$. In particular,
this proposition neither substitutes a bounded-degree statement
for that missing geometric argument nor promotes F13.5 or the
derived-length-three case to Proved.
